\documentclass{beamer}
\usepackage[utf8]{inputenc}
\usepackage{graphicx, epsfig}
\usepackage{amsmath,mathrsfs,amsfonts,amssymb}
\usepackage{floatflt}
\usepackage{epic,ecltree}
\usepackage{mathtext}
\usepackage{fancybox}
\usepackage{fancyhdr}
\usepackage{multirow}
\usepackage{enumerate}
\usepackage{epstopdf}
\usepackage{multicol}
\usepackage{algorithm}
\usepackage[noend]{algorithmic}
\usepackage{tikz}
\usepackage{blindtext}
\usetheme{default}%{Singapore}%{Warsaw}%{Warsaw}%{Darmstadt}
\usecolortheme{default}

\setbeamerfont{title}{size=\Huge}
\setbeamertemplate{footline}[page number]{}

\setbeamertemplate{section in toc}[sections numbered]


\makeatletter
\newcommand\HUGE{\@setfontsize\Huge{35}{40}}
\makeatother    

\setbeamerfont{title}{size=\HUGE}
\beamertemplatenavigationsymbolsempty

% latin bold lower
\newcommand{\ba}{\mathbf{a}} 
\newcommand{\bc}{\mathbf{c}} 
\newcommand{\be}{\mathbf{e}} 
\newcommand{\bff}{\mathbf{f}} % \bf - for bold type
\newcommand{\bg}{\mathbf{g}} 
\newcommand{\bh}{\mathbf{h}} 
\newcommand{\bp}{\mathbf{p}} 
\newcommand{\bq}{\mathbf{q}} 
\newcommand{\bt}{\mathbf{t}} 
\newcommand{\bs}{\mathbf{s}} 
\newcommand{\bu}{\mathbf{u}} 
\newcommand{\bv}{\mathbf{v}} 
\newcommand{\bw}{\mathbf{w}} 
\newcommand{\bx}{\mathbf{x}} 
\newcommand{\by}{\mathbf{y}} 
\newcommand{\bz}{\mathbf{z}} 

% latin bold upper
\newcommand{\bA}{\mathbf{A}} 
\newcommand{\bB}{\mathbf{B}} 
\newcommand{\bC}{\mathbf{C}} 
\newcommand{\bG}{\mathbf{G}} 
\newcommand{\bI}{\mathbf{I}} 
\newcommand{\bJ}{\mathbf{J}} 
\newcommand{\bL}{\mathbf{L}} 
\newcommand{\bM}{\mathbf{M}} 
\newcommand{\bP}{\mathbf{P}}
\newcommand{\bQ}{\mathbf{Q}} 
\newcommand{\bR}{\mathbf{R}} 
\newcommand{\bT}{\mathbf{T}} 
\newcommand{\bU}{\mathbf{U}} 
\newcommand{\bV}{\mathbf{V}} 
\newcommand{\bW}{\mathbf{W}} 
\newcommand{\bX}{\mathbf{X}} 
\newcommand{\bY}{\mathbf{Y}} 
\newcommand{\bZ}{\mathbf{Z}} 

% latin cal upper
\newcommand{\cF}{\mathcal{F}} 
\newcommand{\cG}{\mathcal{G}} 
\newcommand{\cI}{\mathcal{I}} 
\newcommand{\cL}{\mathcal{L}} 
\newcommand{\cM}{\mathcal{M}} 
\newcommand{\cN}{\mathcal{N}} 
\newcommand{\cP}{\mathcal{P}} 
\newcommand{\cS}{\mathcal{S}} 
\newcommand{\cT}{\mathcal{T}} 
\newcommand{\cW}{\mathcal{W}} 
\newcommand{\cX}{\mathcal{X}} 
\newcommand{\cZ}{\mathcal{Z}} 

% latin bb upper
\newcommand{\bbE}{\mathbb{E}} 
\newcommand{\bbI}{\mathbb{I}} 
\newcommand{\bbP}{\mathbb{P}} 
\newcommand{\bbR}{\mathbb{R}} 

% greek bold lower
\newcommand{\bepsilon}{\boldsymbol{\epsilon}} 
\newcommand{\btheta}{\boldsymbol{\theta}} 
\newcommand{\blambda}{\boldsymbol{\lambda}} 
\newcommand{\bpi}{\boldsymbol{\pi}} 
\newcommand{\bmu}{\boldsymbol{\mu}} 
\newcommand{\bsigma}{\boldsymbol{\sigma}} 
\newcommand{\bphi}{\boldsymbol{\phi}} 

% greek bold upper
\newcommand{\bSigma}{\boldsymbol{\Sigma}} 

\DeclareMathOperator*{\argmin}{arg\,min}
\DeclareMathOperator*{\argmax}{arg\,max}
\newcommand{\createdgmtitle}[1]{\title[\hbox to 56mm{Deep Generative Models  \hfill\insertframenumber\,/\,\inserttotalframenumber}]
	{\vspace{1cm} \\ \textbf{Deep Generative Models} \\ {\Huge Lecture #1}}
	\author{Roman Isachenko}
		\institute{\includegraphics[width=0.7cm]{../utils/aimasterslogo} \LARGE{AI Masters}}
	\date{2025, Spring}
}

\usepackage{tikz}
\usetikzlibrary{arrows,shapes,positioning,shadows,trees}

\newcommand\myfootnote[1]{%
  \tikz[remember picture,overlay]
  \draw (current page.south west) +(1in + \oddsidemargin,0.5em)
  node[anchor=south west,inner sep=0pt]{\parbox{\textwidth}{%
      \rlap{\rule{10em}{0.4pt}}\raggedright\scriptsize \textit{#1}}};}

\newcommand\myfootnotewithlink[2]{%
  \tikz[remember picture,overlay]
  \draw (current page.south west) +(1in + \oddsidemargin,0.5em)
  node[anchor=south west,inner sep=0pt]{\parbox{\textwidth}{%
      \rlap{\rule{10em}{0.4pt}}\raggedright\scriptsize\href{#1}{\textit{#2}}}};}
      
\AtBeginSection[]
      {
      	\begin{frame}{Outline}
      		\tableofcontents[currentsection]
      	\end{frame}
      }
      \AtBeginSubsection[]{
      	\begin{frame}{Outline}
      		\tableofcontents[currentsection,currentsubsection]
      	\end{frame}
}
\createdgmtitle{12}

\usepackage{tikz}

\usetikzlibrary{arrows,shapes,positioning,shadows,trees}
%--------------------------------------------------------------------------------
\begin{document}
%--------------------------------------------------------------------------------
\begin{frame}[noframenumbering,plain]
%\thispagestyle{empty}
\titlepage
\end{frame}
%=======
\begin{frame}{Recap of previous lecture}
	\begin{block}{Theorem (continuity equation)}
		If $\bff$ is uniformly Lipschitz continuous in $\bx$ and continuous in $t$, then
		\[
			\frac{d \log p_t(\bx(t))}{d t} = - \text{tr} \left( \frac{\partial \bff(\bx(t), t)}{\partial \bx(t)} \right)
		\]
		\vspace{-0.3cm}
		\[
			\log p_1(\bx(1)) = \log p_0(\bx(0)) - \int_{0}^{1} \text{tr}  \left( \frac{\partial \bff(\bx(t), t)}{\partial \bx(t)} \right) dt.
		\]
		\vspace{-0.3cm}
	\end{block}
	\begin{itemize}
		\item \textbf{Discrete-in-time NF}: evaluation of determinant of the Jacobian costs $O(m^3)$ (we need invertible $\bff$).
		\item \textbf{Continuous-in-time NF}: getting the trace of the Jacobian costs $O(m^2)$ (we need smooth $\bff$).
	\end{itemize}
	\begin{block}{Hutchinson's trace estimator}
		\vspace{-0.3cm}
		\[
	  		\log p_1(\bx(1)) = \log p_0(\bx(0)) - \mathbb{E}_{p(\bepsilon)} \int_{0}^{1} \left[ {\color{violet}\bepsilon^T \frac{\partial \bff}{\partial \bx}} \bepsilon \right] dt.
		\]
	\end{block}
	\myfootnotewithlink{https://arxiv.org/abs/1810.01367}{Grathwohl W. et al. FFJORD: Free-form Continuous Dynamics for Scalable Reversible Generative Models, 2018} 
\end{frame}
%=======
\begin{frame}{Recap of previous lecture}
	\begin{block}{Forward pass (Loss function)}
		\vspace{-0.7cm}
		\[
			L(\bx) = - \log p_1(\bx(1) | \btheta) = - \log p_0(\bx(0)) + \int_{0}^{1} \text{tr} \left( \frac{\partial \bff_{\btheta}(\bx(t), t)}{\partial \bx(t)} \right) dt
		\]
	\end{block}
	\vspace{-0.5cm}
	\begin{block}{Adjoint functions}
		\vspace{-0.3cm}
		\[
			\ba_{\bx}(t) = \frac{\partial L}{\partial \bx(t)}; \quad \ba_{\btheta}(t) = \frac{\partial L}{\partial \btheta(t)}.
		\]
		\vspace{-0.5cm}
	\end{block}
	\begin{block}{Theorem (Pontryagin)}
		\vspace{-0.5cm}
		\[
		\frac{d \ba_{\bx}(t)}{dt} = - \ba_{\bx}(t)^T \cdot \frac{\partial \bff_{\btheta}(\bx(t), t)}{\partial \bx}; \quad \frac{d \ba_{\btheta}(t)}{dt} = - \ba_{\bx}(t)^T \cdot \frac{\partial \bff_{\btheta}(\bx(t),  t)}{\partial \btheta}.
		\]
		\vspace{-0.7cm}
		\begin{align*}
			\frac{\partial L}{\partial \btheta(0)} &= \ba_{\btheta}(0) =  - \int_{1}^{0} \ba_{\bx}(t)^T \frac{\partial \bff_{\btheta}(\bx(t), t)}{\partial \btheta(t)} dt + 0\\
			\frac{\partial L}{\partial \bx(0)} &= \ba_{\bx}(0) =  - \int_{1}^{0} \ba_{\bx}(t)^T \frac{\partial \bff_{\btheta}(\bx(t), t)}{\partial \bx(t)} dt + \frac{\partial L}{\partial \bx(1)}
		\end{align*}
	\end{block}
	\myfootnotewithlink{https://arxiv.org/abs/1806.07366}{Chen R. T. Q. et al. Neural Ordinary Differential Equations, 2018}   
\end{frame}
%=======
\begin{frame}{Recap of previous lecture}
	\begin{block}{Forward pass}
		\vspace{-0.3cm}
		\[
		\bx(1) = \bx(0) + \int^{1}_{0} \bff_{\btheta}(\bx(t), t) d t \quad \Rightarrow \quad \text{ODE Solver}
		\]
		\vspace{-0.4cm}
	\end{block}
	\begin{block}{Backward pass}
		\vspace{-0.5cm}
		\begin{equation*}
			\left.
			{\footnotesize 
				\begin{aligned}
					\frac{\partial L}{\partial \btheta(0)} &= \ba_{\btheta}(0) =  - \int_{1}^{0} \ba_{\bx}(t)^T \frac{\partial \bff_{\btheta}(\bx(t), t)}{\partial \btheta(t)} dt + 0 \\
					\frac{\partial L}{\partial \bx(0)} &= \ba_{\bx}(0) =  - \int_{1}^{0} \ba_{\bx}(t)^T \frac{\partial \bff_{\btheta}(\bx(t), t)}{\partial \bx(t)} dt + \frac{\partial L}{\partial \bx(1)} \\
					\bx(0) &= - \int^{1}_{0} \bff_{\btheta}(\bx(t), t) d t  + \bx(1).
				\end{aligned}
			}
			\right\rbrace
			\Rightarrow
			\text{ODE Solver}
		\end{equation*}
		\vspace{-0.4cm} 
	\end{block}
	\textbf{Note:} These scary formulas are the standard backprop in the discrete case.
	\myfootnotewithlink{https://arxiv.org/abs/1806.07366}{Chen R. T. Q. et al. Neural Ordinary Differential Equations, 2018}   
\end{frame}
%=======
\begin{frame}{Outline}
	\tableofcontents
\end{frame}
%=======
\begin{frame}{Summary}
	\begin{itemize}
		\item 
	\end{itemize}
\end{frame}
\end{document} 